\section{Lepton-proton scattering}\label{LPScatt}


\begin{figure}[b]
\begin{center}
\includegraphics[scale=0.45,angle=0]{Jaxo/elasticLP.pdf}
\qquad
\includegraphics[scale=0.45,angle=0]{Jaxo/OPE.pdf}
\end{center}
\caption{
\label{fig:rc:elLP}
Generic $\ell p$ scattering (left panel), tree-level Born diagram (right panel).}
\end{figure}
%--


Scattering experiments allow experimenters to probe the structure of the target. From a measurement of the unpolarized $\ell p$ scattering cross section, one can for instance determine the proton electric and magnetic Sachs form factors (FFs), $G_E(Q^2)$ and $G_M(Q^2)$, functions of the momentum transfer $t=-Q^2$. 
The Born cross section of elastic $\ell p$ scattering, given by the tree-level one-photon exchange (OPE) diagram in Fig.~\ref{fig:rc:elLP} (right panel), can be written in a compact form as

\beq
\eqlab{BornXS}
\frac{\mathrm{d} \sigma^{(0)}}{\mathrm{d} \Omega} = \frac{\sigma_\text{Mott} \,\tau}{\varepsilon(1+\tau)}\,\sigma_R,
\eeq
with the Mott cross section for scattering off a point-like particle
\beq \label{SigmaRed}
\sigma_\text{Mott}=\frac{\alpha^2 \epsilon_2 \cos^2{\theta/2}}{4 \epsilon_1^3 \sin^4{\theta/2}},
\eeq
the reduced Born cross section
\beq
\sigma_R =G_M^2(Q^2)+\frac{\varepsilon}{\tau}\, G_E^2(Q^2),\eqlab{reducedCS}
\eeq
and the 
photon polarization parameter 
\beq 
\label{eq:epsilon_def}
\varepsilon = \frac{\nu^2 - M^4 \tau ( 1 + \tau )}{ \nu^2 + M^4 \tau ( 1 + \tau )  ( 1 - 2\varepsilon_{\ell})}.
\eeq
 Here, we introduced the 
 dimensionless quantities $\tau=Q^2/4M^2$, $\varepsilon_\ell = 1/(2\tau_\ell)$ and $\tau_\ell = Q^2 / (4 m_\ell^2)$, and the crossing-symmetric variable $ \nu \equiv k_1 \cdot (p_1 + p_2)/2 = (s - u)/4$, with $M$ ($m_\ell)$ the proton (lepton) mass, and where $s$ and $u$ denote the usual Mandelstam variables of the elastic scattering process of Fig.~\ref{fig:rc:elLP}. Furthermore, $\alpha \simeq 1/137$ is the usual fine structure constant, $\epsilon_1$ ($\epsilon_2$) is the energy of the incoming (outgoing) lepton, and $\theta$ is the scattering angle in the laboratory frame, see \Eqref{eq2}.

The extraction of the FFs, cf.\ \Eqref{reducedCS}, is done via Rosenbluth separation  \cite{Rosenbluth:1950yq}, which  requires measurements at the same momentum transfer $Q^2$ but for different energies and angles. To recover the Born cross section, the experimentally measured cross section needs to be corrected for radiative corrections, cf. Refs.~\cite{Maximon:2000hm,Mo:1968cg,Ent:2001hm,Gramolin:2014pva,Gerasimov:2015aoa,Bystritskiy:2007hw}. These change not only the absolute value of the cross section but also its dependence on the relevant kinematical variables. In general, the  $\mathcal{O}(\al^2)$ Born cross section has to be corrected by additional virtual photons or inelastic scattering with emission of real bremsstrahlung. 



 
\subsection{Radiative corrections}\label{Sec:RadCorrAdrian}

In this subsection, we discuss the calculation of cross sections in a
perturbative expansion of the electromagnetic coupling $\alpha$. While
many remarks are valid for arbitrary processes, we will illustrate the
general procedure for elastic $\ell p$ scattering \beq
\ell( k_1 , h ) + p( p_1, \lambda ) \to \ell( k_2, h^\prime) + p(p_2, \lambda^\prime) \label{eq:eq1}
\eeq
where $k_1(k_2)$ and $p_1(p_2)$ are the incoming (outgoing) lepton and proton momenta, and $h(h^\prime)$ and $\lambda(\lambda^\prime)$ the respective helicities.

The Born or leading order (LO) cross section for elastic $\ell$p scattering, given in \Eqref{BornXS}, is of
$\mathcal{O}(\alpha^2 Z^2)$ and is a function of $G_E(Q^2)$ and $G_M(Q^2)$.  These form factors are depicted as grey blobs
in~\Figref{rc:LO} and, for each appearance, a factor $Z$, the total charge, is included in
the counting of the couplings. To improve the theory description,
next-to-leading order (NLO) corrections to \Eqref{BornXS} have been
computed in Refs.~\cite{Maximon:2000hm, Mo:1968cg, Ent:2001hm,
  Gramolin:2014pva, Gerasimov:2015aoa, Bystritskiy:2007hw}. The NLO
corrections can be decomposed into several gauge-invariant parts. The one-photon exchange (OPE) contributions are of the form
$\mathcal{O}(\alpha^3 Z^2)$ and also include vacuum polarisation (VP)
corrections. They are illustrated in Figs.~\ref{fig:rc:lNLO} and
\ref{fig:rc:vpNLO}, respectively. The TPE
contributions, depicted in
Figs.~\ref{fig:rc:pNLO} and \ref{fig:rc:messNLO}, are $\mathcal{O}(\alpha^3 Z^3)$. Finally, corrections to the
proton line are of $\mathcal{O}(\alpha^3 Z^4)$, with a sample
contribution shown in \Figref{rc:FF-NLO}. These corrections change
not only the absolute value of the observables but also their
dependence on the relevant kinematical variables, e.g., the square of
the transferred momentum $Q^2$ and the photon polarization
parameter $\varepsilon$, defined in~\Eqref{epsilon_def}.

If the radiative corrections are small, an NLO calculation might be
sufficient. In particularly simple situations, these corrections can be
included as a multiplicative factor. The
error associated to this procedure is at the percent level. Radiative corrections can lead to
contributions that are enhanced by large logarithms of the form
$L=\ln{(Q^2/m_\ell^2)}$.  As we will discuss below, these are associated with hard
collinear radiation. Another source of large logarithms are stringent
experimental cuts that restrict the phase space for soft emission. As a consequence, modern high-precision experiments are required to introduce higher order corrections
and a more sophisticated treatment of real radiation.



\begin{figure}[ht]
\begin{center}
\includegraphics[scale=0.45,angle=0]{Jaxo/LO.pdf}
\end{center}
\caption[]{
\label{fig:rc:LO}
Schematic illustration of the squared LO amplitude
$\AMP^{(0)}\AMP^{(0)\,*}$ for lepton-proton scattering. Leptons
(protons) are depicted as single (double) lines and the grey blobs
represent the form factors.  }
\end{figure}
%---------------------------------------

%---------------------------------------
\begin{figure}[ht]
\begin{center}
\includegraphics[scale=0.45,angle=0]{Jaxo/lNLO.pdf}
\end{center}
\caption[]{
\label{fig:rc:lNLO}
An example contribution $\mathcal{O}(\alpha^3\,Z^2)$ to the squared NLO
amplitude with photon emission restricted to the lepton line. The
virtual (v) cut represents $\AMP^{(1)}\,\AMP^{(0)\,*}$ whereas the
real (r) cut represents $|\AMP_\gamma^{(0)}|^2$. }
\end{figure}
%---------------------------------------

%---------------------------------------
\begin{figure}[ht]
\begin{center}
\includegraphics[scale=0.45,angle=0]{Jaxo/vpNLO.pdf}
\end{center}
\caption[]{
\label{fig:rc:vpNLO}
An NLO correction due to VP contributions. The virtual (v) cut
represents $\AMP^{(1)}\,\AMP^{(0)\,*}$. }
\end{figure}
%---------------------------------------

%---------------------------------------
\begin{figure}[ht]
\begin{center}
\includegraphics[scale=0.45,angle=0]{Jaxo/pNLO.pdf}
\end{center}
\caption[]{
\label{fig:rc:pNLO}
An example contribution to the squared NLO amplitude
$\mathcal{O}(\alpha^3 Z^3)$ with elastic TPE. The
virtual (v) cut represents $\AMP^{(1)}\,\AMP^{(0)\,*}$ whereas the
real (r) cut represents $|\AMP_\gamma^{(0)}|^2$. }
\end{figure}
%---------------------------------------

%---------------------------------------
\begin{figure}[ht]
\begin{center}
\includegraphics[scale=0.45,angle=0]{Jaxo/messNLO.pdf}
\end{center}
\caption[]{
\label{fig:rc:messNLO}
An example contribution to the squared NLO amplitude with inelastic
TPE. The virtual cut (v) represents
$\AMP^{(1)}\,\AMP^{(0)\,*}$ whereas the real cut (r) represents
$|\AMP_\gamma^{(0)}|^2$. }
\end{figure}
%---------------------------------------
%---------------------------------------
%\restylefloat{figure}
\begin{figure}[ht]
\begin{center}
\includegraphics[scale=0.45,angle=0]{Jaxo/FF-NLO.pdf}
\end{center}
\caption{
\label{fig:rc:FF-NLO}
An example contribution $\mathcal{O}(\alpha^3\,Z^4)$ to the squared NLO
amplitude with photon emission restricted to the proton line. The
virtual (v) cut represents $\AMP^{(1)}\,\AMP^{(0)\,*}$ whereas the
real (r) cut represents $|\AMP_\gamma^{(0)}|^2$. }
\end{figure}
%---------------------------------------

This leads us to consider fully differential computations of cross
sections which allows one to obtain precise predictions for fiducial cross
sections. To keep full generality, the phase-space integration needs
to be adaptable to the experimental setting. As a result, it is
necessary to revert to numerical integrations, which are most
efficiently carried out using Monte Carlo (MC) techniques.

For the LO cross section $\dd \sigma^{(0)}\sim\alpha^2\,Z^2$, the
tree-level matrix element squared $|\AMP^{(0)}|^2$, depicted in
\Figref{rc:LO}, is to be integrated over the two-body phase space
$\Phi_2$. Thus,
\begin{equation}\label{mcLO}
\dd\sigma^{(0)} = \int \dd\Phi_2\, |\AMP^{(0)}|^2\, S(k_2,p_2)\,,
\end{equation}
where the so-called measurement function $S(k_2,p_2)$ defines the
quantity to be computed in terms of the momenta of the final state
particles. For simplicity, trivial elements such as the flux 
factor are also included in $S$. If a cut is to be
applied, $S$ is set to 0 if the event does not pass the cut. 

Going beyond LO, we encounter ultraviolet (UV) and infrared (IR)
divergences. After renormalization, in a QED calculation where all
fermion masses are kept at their finite values, only soft IR
divergences remain. These soft singularities cancel between the
virtual and real corrections. At NLO this is illustrated in
\Figref{rc:lNLO}, where some corrections due to photon emission from
the lepton line are shown. The virtual corrections are obtained as
\begin{equation}\label{mcVNLO}
\dd\sigma^{(1)}_{\rm v} = \int \dd\Phi_2\, 2\re(\AMP^{(1)}\AMP^{(0)\,*})\,
S(k_2,p_2)\, ,
\end{equation}
whereas the real corrections are given by 
\begin{equation}\label{eq:mcRNLO}
\dd\sigma^{(1)}_{\rm r} = \int \dd\Phi_3\, |\AMP_\gamma^{(0)}|^2\,
S(k_2,p_2,k_\gamma)\, .
\end{equation}
Both terms are $\mathcal{O}(\alpha^3\,Z^2)$. Here $\AMP^{(1)}$ is the one-loop amplitude of the process \Eqref{eq1} and  $\AMP_\gamma^{(0)}$
is the tree-level amplitude with an additional photon of momentum
$k_\gamma$ in the final state. The integration in \Eqref{mcRNLO} has
to be carried out over the three-particle final state $\Phi_3$. Also,
the definition of $q$ is more subtle now, since $k_1-k_2\neq p_2-p_1$.





To give mathematical meaning individually to $\dd \sigma^{(1)}_{\rm v}$
and $\dd \sigma^{(1)}_{\rm r}$, the IR singularities need to be
regularized. Many calculations in QED still use a (small) fictitious
photon mass to do so. Then, the IR singularities manifest themselves as
logarithms of this mass. Another approach that is more in line with
the tremendous progress in computational techniques that has been made
for QCD calculations is to use dimensional regularization (i.e. work in
$4-2\epsilon$ dimensions) also for IR singularities. In this case,
IR singularites show up as poles $1/\epsilon$. See \secref{IRsubtraction} for a detailed discussion. Independent of the
chosen regularization, in the combination of $\dd \sigma^{(1)}_{\rm v}$
and $\dd \sigma^{(1)}_{\rm r}$ the remnants of the IR singularities
(i.e. the photon mass logarithms or the $1/\epsilon$ poles)
cancel. This is ensured by the well-known limiting behavior of the
real matrix element in the soft limit ($\mathcal{S}@$)~\cite{Yennie:1961ad}, where it
can be written as an eikonal factor $\mathcal{E}$ times the Born
matrix element
\begin{equation}\label{eikonal}
\mathcal{S}@ |\AMP_\gamma^{(0)}|^2 = \mathcal{E}\,|\AMP^{(0)}|^2 . 
\end{equation}
Hence, in the complete NLO corrections to Eq.~\ref{mcLO},
\begin{equation}\label{mcNLO}
\dd \sigma^{(1)} = \dd \sigma^{(1)}_{\rm v} + \dd \sigma^{(1)}_{\rm r} ,
\end{equation}
the regularization can be removed (i.e. the photon mass
or $\epsilon$ set to zero). For this cancellation to occur, the
observable has to be IR safe. This means that the observable must not
change whether or not an arbitrarily soft photon is emitted. At a
technical level, we must require
\begin{equation}\label{irsafe}
  \lim_{k_\gamma \to 0} S(k_2,p_2,k_\gamma) = S(k_2,p_2) .
\end{equation}  
From a mathematical point of view, it is sufficient if Eq.~\ref{irsafe}
holds in the strict limit. However, if experimental cuts (directly or
indirectly affecting real radiation) induce a difference between
$S(k_2,p_2,k_\gamma)$ and $S(k_2,p_2)$ for finite, but small photon
energies $\Delta{E_\gamma}$, there can be a large logarithm
$\ln\left(\Delta{E_\gamma^2}/Q^2\right)$ as a remnant. Such (soft) logarithms can lead
to enhanced QED corrections.

Another source of large corrections is collinear emission. While
collinear singularites are regulated by the fermion masses, they still
can also lead to enhanced logarithmic terms $\alpha L$ after
combination of real and virtual corrections. These logarithms often
form the dominant part of the higher-order corrections. This also
entails that in fully differential QED calculations -- contrary to QCD
calculations -- it is not possible to set the lepton masses to
zero. Furthermore, the neglect of hard collinear emission potentially leads
to a loss of accuracy. Thus, using the soft approximation for the real
matrix elements can have severe implications on the accuracy of the
results. 

The extraction of the IR poles from \Eqref{mcRNLO} for an arbitrary
IR-safe observable, i.e. an arbitrary $S(k_2,p_2,k_\gamma)$ satisfying
Eq.~\ref{irsafe}, is by now a standard procedure at NLO and the focus
turns towards next-to-next-to leading order (NNLO) applications as discussed in Ref.~\cite{TorresBobadilla:2020ekr}. Again,
at NLO there are two widely used options, the slicing method and the
subtraction method.

In the slicing method, the phase space is split into two parts,
depending on the photon energy
\begin{align}\label{slicing}
  \dd\sigma^{(1)}_{\rm r}\big|_\text{sl}
  &=
  \int\displaylimits_{E_\gamma<\delta} \hspace*{-5pt}
  \dd\Phi_3\, |\AMP_\gamma^{(0)}|^2\,S(k_2,p_2)
\nonumber \\
 &+ \int\displaylimits_{E_\gamma\ge\delta} \hspace*{-5pt}
 \dd\Phi_3\, |\AMP_\gamma^{(0)}|^2\,S(k_2,p_2,k_\gamma)
\, .
\end{align}
In the part, where the photon energy is larger than a chosen resolution
parameter $\delta$, the integration can be carried out for a massless
photon without encountering singularities. In the part, where the
photon energy is smaller than the resolution parameter, the soft
(eikonal) approximation Eq.~\ref{eikonal} is used for the matrix
element. The integration over the photon phase space then simplifies
and can be carried out analytically. In this term, the photon mass
logarithms that cancel those of the virtual
corrections are generated. This method relies on a suitable choice for $\delta$. It should not be too large, to ensure
that the soft approximation is good enough. It should also not be too
small, otherwise numerical problems (large cancellations between the
resolved and unresolved parts) occur. For more details, see Sec.~\ref{sec:Implementations}.

In the subtraction method a term (the soft limit of the integrand) is
subtracted and added back
\begin{align}\label{subtraction}
  &\dd\sigma^{(1)}_{\rm r}\big|_\text{su}  =  \int \dd\Phi_3\ \mathcal{S}@ |\AMP_\gamma^{(0)}|^2  \,S(k_2,p_2) \ +\nonumber \\
   &\int \dd\Phi_3 \bigg(|\AMP_\gamma^{(0)}|^2S(k_2,p_2,k_\gamma) -\mathcal{S}@ |\AMP_\gamma^{(0)}|^2S(k_2,p_2)\bigg).
\end{align}
The first term on the r.h.s. of Eq.~\ref{subtraction} is very similar to
the corresponding term in Eq.~\ref{slicing}. The only differences are
that the photon is treated as massless and the integration is done
over the full phase space in Eq.~\ref{subtraction}. This results in
IR $1/\epsilon$ poles that cancel against the IR poles of the virtual
contributions. The second term on the r.h.s. of Eq.~\ref{subtraction} is
finite and can be integrated in four dimensions. An efficient numerical
evaluation of this term including the subtraction requires an adapted
phase-space parameterization. Within the subtraction method, there is
no motivation to split the real corrections into hard and soft
parts. It is actually simpler to always include the full real matrix
element and make no approximation whatsoever.

The OPE corrections due to emission from the lepton line
$\mathcal{O}(\alpha^3\,Z^2)$ and the VP corrections, illustrated in Figs.~\ref{fig:rc:lNLO} and
\ref{fig:rc:vpNLO}, are simple from a conceptual
point of view as they involve standard QED pointlike interactions
only. However, there are also corrections of $\mathcal{O}(\alpha^3\,Z^3)$
that involve multiple exchange of photons between the two fermion
 lines, in particular the notorious TPE
corrections. The class of corrections depicted in \Figref{rc:pNLO} is
more complicated than those of \Figref{rc:lNLO} for several
reasons. Even if the proton was pointlike, these corrections now
involve one-loop box diagrams at NLO. In addition, there is the complication that the
photon-proton vertex is more involved than for a pointlike
particle. Furthermore, for $\ell p$ scattering there are also
contributions where the intermediate state is not just a proton, as
depicted in \Figref{rc:messNLO}. These aspects will be discussed in
\secref{TPE}.



A final class of NLO corrections are those where the emission of
additional photons is restricted to the proton line. They are
$\mathcal{O}(\alpha^3\,Z^4)$ and a sample contribution is depicted in
\Figref{rc:FF-NLO}. It is tempting to include these diagrams into
the definition of the form factors. However, they contain IR
divergences that, once more, cancel between the real and virtual
parts. Thus, from a practical point of view it is more convenient to
explicitly include these diagrams into the radiative corrections and define the form factors
accordingly, i.e. exclude all QED effects from their definition. 
The numerical impact of these contributions is rather minor, as they
do not generate collinear logarithms. Hence, it is often possible to
simply neglect the $\mathcal{O}(\alpha^3\,Z^4)$ contribution. 





As discussed above, at each order in $\alpha$ the QED corrections can
be enhanced by a collinear logarithm $L$ and possibly a soft logarithm
$\ln\left(\Delta{E_\gamma}^2/Q^2\right)$. Hence, the corrections can be considerably
larger than naive scaling with $\alpha$ implies. Thus, for very
precise predictions it is necessary to go beyond NLO. Regarding
OPE corrections restricted to emission from the lepton line
$\mathcal{O}(\alpha^4\,Z^2)$, the computations have been extended to
NNLO. This involves double-virtual, real-virtual, and double-real
corrections, as illustrated in \Figref{rc:NNLO}. Results have been
obtained using photon-mass regularization and the slicing
method in Refs.~\cite{Bucoveanu:2018soy,CarloniCalame:2020yoz} as well as using
dimensional regularization with the subtraction
method in Ref.~\cite{Banerjee:2020rww}. This class of corrections is a gauge
invariant subset of the full corrections. The virtual part involves
only the vector part of the two-loop heavy-lepton form
factor, cf. Refs.~\cite{Mastrolia:2003yz,Bonciani:2003ai,Bernreuther:2004ih}. For
lepton masses much smaller than the proton mass, these corrections are
typically also dominant. They involve up to two additional photon
emissions and a precise description how these photons are treated is
required. This amounts to a definition of
$S(k_2,p_2,k_{\gamma_1},k_{\gamma_2})$. In particular, both photons
can be collinear, leading to enhanced NNLO correction of the form
$(\alpha\,L)^2$. 

%---------------------------------------
%\restylefloat{figure}
\begin{figure}[ht]
\begin{center}
\includegraphics[scale=0.45,angle=0]{Jaxo/NNLO.pdf}
\end{center}
\caption[]{
\label{fig:rc:NNLO}
Example of a NNLO contribution $\mathcal{O}(\alpha^4 Z^2)$ to the
squared amplitude for $\ell p$ scattering. The double-virtual
(vv) cut represents $\AMP^{(2)}\AMP^{(0)\,*}$, the real-virtual (rv)
cuts $\AMP^{(0)}_\gamma\AMP^{(1)\,*}_\gamma$ and
$\AMP^{(1)}_\gamma\AMP^{(0)\,*}_\gamma$, and the double-real (rr) cut
$|\AMP^{(0)}_{\gamma\gamma}|^2$.}
\end{figure}
%---------------------------------------

Corrections at NNLO that go beyond the OPE approximation lead to
substantially more complicated calculations. If the proton is treated
as pointlike, the NNLO results for electron-muon
scattering in Refs.~\cite{Engel:2022kde, Broggio:2022htr} can be adapted by
simple replacements $m_\mu\to{M}$ and $m_e\to{m}_\ell$. This requires the computation of the two-loop amplitude $\AMP^{(2)} $(for double-virtual) involving (crossed) double-box diagrams~\cite{Bonciani:2021okt}. For the real-virtual corrections $\AMP_\gamma^{(1)}$ is required where one-loop pentagon diagrams contribute. Efficient techniques for one-loop computations have been developed in connection with high multiplicity processes at hadron colliders, as reviewed in Ref.~\cite{Ellis:2011cr}.
As part of this topical collection, the effect of pointlike NNLO corrections were compared to present uncertainties from hadronic correction~\cite{McMule}.
It turns out that the foremer can be sizeable and need to be properly taken into account.
Beyond pointlike photon-proton interaction make the formidable task of the
two-loop amplitude computation even more
daunting. At NNLO, depending on the experimental set up, it might also
be necessary to include the process with  emission of an additional
lepton pair which also has been computed for electron-muon
scattering in Ref.~\cite{Budassi:2021twh}.

First steps to move beyond NNLO have already been taken. The form factor $\gamma^*\to\ell^+\ell^-$ is known
at three loops~\cite{Fael:2022rgm} and the subtraction method for
massive QED has been extended beyond NNLO~\cite{Engel:2019nfw}. Hence,
a next-to-next-to-next-to leading order (NNNLO) calculation of OPE contributions $\mathcal{O}(\alpha^5\,Z^2)$
seems to be feasible within a reasonable time frame. However, a full
computation of $2\to{2}$ scattering processes with massive particles at
NNNLO is currently out of reach.

As alluded to above, the corrections are often dominated by large
logarithms associated with collinear or soft emission. A convenient
method resumming all contributions of large logarithms at all orders
has been proposed in Ref.~\cite{Kuraev:1985hb}. All leading contributions
of order $(\alpha L)^n$ are taken into account.
%resulting in permille precision.
The accuracy can be further improved calculating non-leading
contributions of order $\alpha (\alpha L)^n$ as a $K$-factor. The application to $ep$ elastic scattering can be found in
Ref.~\cite{Bystritskiy:2006ju}, and, including hard photon emission,
in Ref.~\cite{Kuraev:2013dra}. 

A more generic approach to include logarithmically enhanced
corrections is to use QED parton showers. A recent review of this
activity can be found in Ref.~\cite{Frixione:2022ofv} and will be discussed
in some more detail in Sec.~\ref{sec:Implementations}.  Parton
showers allow for the inclusion of the leading logarithmic
contributions in a process independent way. Combined with fixed-order
calculations this is a powerful tool to improve further the precision
for fully differential cross sections. Including next-to-leading
logarithms from initial-state (or final-state) collinear emission is also possible in
a generic way~\cite{Bertone:2019hks}. However, in general the
systematic inclusion of subleading logarithms is observable dependent
and some examples will be discussed in Sec.~\ref{sec:ho}. Contrary
to QCD, in QED it is typically not required to actually resum the
logarithms. The suppression by $\alpha$ is strong enough to ensure
that $\alpha\,L$ is still reasonably small. Hence, the inclusion of
the logarithms at the first few orders in $\alpha$ beyond the
fixed-order results is sufficient.



\subsection{Experimental observables}
In the following, we discuss different experimental observables and how they are affected by the LO in $\alpha$ virtual corrections, cf.\ Eq.~\ref{mcVNLO}. Of particular interest are the unpolarized cross section, the longitudinal and transverse polarization transfer observables $P_t$ and $P_l$, as well as the target and beam asymmetries $A_n$ and $B_n$, introduced in Secs.~\ref{sec:PTobservable} and \ref{sec:beam_targ_normal} respectively.

\subsubsection{Elastic lepton-proton scattering}

Let us start by presenting the tensor decomposition of elastic $\ell p$ scattering. 
Taking into account discrete symmetries, assuming parity and time reversal invariance of the electromagnetic interaction, this process can be described in terms of $6$ independent complex scalar amplitudes: $\cG_M$ and $\cF_i$ ($i=2,\dots 6)$, which are generally functions of two
independent kinematical variables, e.g., the Mandelstam variables $t=(k_1-k_2)^2$ and $s=(k_1+p_1)^2$.  The helicity amplitudes of the process can be split into a sum of helicity-conserving and helicity-flip  contributions, see, e.g.,\ Refs.~\cite{Guichon:2003qm,Gorchtein:2004ac,Gakh:2014zva},

\begin{eqnarray}
\label{eq:str_ampl}
\AMP_{h^\prime \lambda^\prime, h \lambda}  &=&  \AMP_{h^\prime \lambda^\prime, h \lambda}^{\mathrm{non-flip}}  +  \AMP_{h^\prime \lambda^\prime, h \lambda}^{\mathrm{flip}},\\
\AMP_{h^\prime \lambda^\prime, h \lambda}^{\mathrm{non-flip}} &=&  \frac{4\pi \al}{Q^2}\, \bar{u}(k_2,h^\prime) \gamma_\mu u(k_1,h)\;  \bar{N}(p_2,\lambda^\prime)\nn\\&& \times\left( \cG_M \gamma^\mu -   \cF_2 \frac{P^{\mu}}{M} +  \cF_3 \frac{\gamma \cdot K P^{\mu}}{M^2} \right)\nn\\ &&\times N(p_1,\lambda),  \label{eq:str_ampl1}  \\
 \AMP_{h^\prime \lambda^\prime, h \lambda}^{\mathrm{flip}} &=& \frac{4\pi \al}{Q^2} \frac{m_\ell}{M}\, \bar{u}(k_2,h^\prime) u(k_1,h) \; \bar{N}(p_2,\lambda^\prime)\nn\\ &&\times\left( \cF_4  +  \frac{\gamma \cdot K}{M} \cF_5  \right) N(p_1,\lambda)  \nonumber \\
 &+ &\frac{4\pi \al}{Q^2} \frac{m_\ell}{M}\, \bar{u}(k_2,h^\prime) \gamma_5 u(k_1,h) \;\nn\\ &&\times\bar{N}(p_2,\lambda^\prime) \cF_6 \gamma_5 N(p_1,\lambda), \label{eq:str_ampl2}
\end{eqnarray}

where $q=k_1-k_2=p_2-p_1$, $K=(k_1+k_2)/2$,
$P=(p_1+p_2)/2$, and $\gamma \cdot a \equiv \gamma^\mu a_\mu$.  
The $6$ complex amplitudes $\cG_M$ and $\cF_i$, sometimes called generalized FFs, fully describe the
spin structure of reaction Eq.~\ref{eq:eq1} for any
number of exchanged virtual photons. In the limit of $m_\ell \to 0$, the contribution from the helicity-flip amplitudes to observables vanishes, see Eq.~\ref{eq:str_ampl2}. Relations between the helicity amplitudes and the generalized FFs can be found f.i.\ in Ref.~\cite{Tomalak:2014dja}.

In the Born approximation, i.e., considering only the tree-level OPE diagram with proton FFs in \Figref{rc:elLP} (right panel), the number of amplitudes in \Eqref{str_ampl} reduces to $2$. To be more precise, the Born amplitudes read
\begin{eqnarray}
\cG_M^\text{Born}(Q^2, s)&=&G_{M}(Q^2), \\
%\cG_E^\text{Born}(q^2, s)&=&G_{E}(q^2), \\
\cF_2^\text{Born}(Q^2,
s)&=&F_{2}(Q^2),  \\
\cF_i^\text{Born}(Q^2, s)&=&0 \quad (i=3,\, \dots, 6),
\label{eq:eq6}
\end{eqnarray}
where $F_{2}(Q^2)$ is the Pauli
 FF of the proton. The fact that the FFs  in the space-like region are 
real functions of the virtuality of the exchanged photon is a consequence of unitarity, the spin-$1$ nature of the
virtual photon, parity conservation, and the identity of the initial
and final states, see Refs.~\cite{Halzen:1984mc,Tomasi-Gustafsson:2011sda}
for an explicit derivation.


To separate the Born contribution from effects of virtual radiative corrections to the elastic scattering, we define the following
decomposition of the scalar amplitudes:
    \begin{eqnarray}
\cG_M(Q^2, s)&=&G_M(Q^2)+\Delta\cG_M(Q^2, s), \\
 \cG_E(Q^2, s)&=&G_E(Q^2)+\Delta\cG_E(Q^2, s),\\
 \cF_i(Q^2, s)&=&\Delta\cF_i(Q^2, s) \quad (i=3,\, \dots, 6),\quad
\label{eq:eq8}
\end{eqnarray}
where $\cG_E$ has been introduced via
\beq
\cF_2(Q^2, s)=\frac{1}{1+\tau}\left[\cG_M(Q^2, s)-\cG_{E}(Q^2, s)\right].
\eeq
The order of magnitude of these quantities is given by
$ G_{M,E}\sim\alpha^0$ and 
$\Delta\cG_M(q^2,s)\sim\Delta\cG_E(q^2, s)~
\sim\Delta\cF_i(q^2, s) (i=3,\, \dots, 6)\sim\alpha  $.

\subsubsection{Laboratory frame}



The proton
and lepton four-momenta in the laboratory frame can be written as
\begin{eqnarray}\label{eq:eq2}
p_1&=&(M,\boldsymbol{0}),\\
p_2&=&(E_2, \boldsymbol{p_2}),\\
k_1&=&(\epsilon_1, \boldsymbol{k_1}),\\ k_2&=&(\epsilon_2, \boldsymbol{k_2}),
\end{eqnarray}
where three momenta are denoted by bold symbols. The scattered lepton
energy is written in terms of $\theta $ as
\begin{eqnarray}    
\epsilon_2&=& \frac{1}
{(\epsilon_1+M)^2-\boldsymbol{k_1}^{\!2}\cos^2\theta }\nonumber\\
&\times& \Bigg[(\epsilon_1+M)(M\epsilon_1 +
m_\ell^2)\nonumber\\
&&+\boldsymbol{k_1}^{\!2}\cos\theta\sqrt{ M^2-m_\ell^2\sin^2\theta}\Bigg],\qquad
\label{eq:eqe12}
\end{eqnarray}
while the Mandelstam variables $s$ and $t$ are expressed as
\begin{eqnarray}\label{eq:eqq2}
t&=&-Q^2=2m_\ell^2-2(\epsilon_1\epsilon_2-\vert\boldsymbol{k_1}\vert \vert\boldsymbol{k_2}\vert\cos\theta),\qquad\\
s&=& M^2 + m_\ell^2 + 2 M \epsilon_1.
\end{eqnarray}
The differential cross section in terms of the  matrix element squared is given by
\begin{eqnarray}
\dd\sigma&=&\frac{(2\pi )^4}{4I}\vert{\AMP}\vert^2\frac{\dd{\boldsymbol {k_2}}\,\dd{\boldsymbol{
p_2}}}{(2\pi )^64\epsilon_2E_2}\nonumber\\
&\times&\,\delta^{(4)}(k_1+p_1-k_2-p_2),
\label{eq:eq10}
\end{eqnarray}
with the invariant $I^2=(k_1\cdot p_1)^2-m_\ell^2M^2$ and the
energy of the recoil proton $E_2$.

For the case where the scattered lepton is detected in the final state, one obtains 
\begin{equation}
\frac{\dd\sigma}{\dd\Omega}=\frac{\vert \AMP \vert^2}{(4\pi)^2 \,4M}\frac{\boldsymbol{k_2}^{\!2}}{\mathfrak{D}\vert\boldsymbol{ k_1}\vert},
%{\frac{\dd\sigma}{\dd\Omega}=\frac{\alpha ^2}{4M}\frac{\boldsymbol{k_2}^{\!2}}{\mathfrak{d}\vert\boldsymbol{ k_1}\vert}\frac{\vert\overline{\cal M}\vert^2}{Q^4},}
\label{eq:eq11}
\end{equation}
where $\mathfrak{D}=(M+\epsilon_1)\vert\boldsymbol{ k_2}\vert-\epsilon_2\vert\boldsymbol{ k_1}\vert\cos\theta $, and $\dd \Omega$ is the differential solid angle of the scattered lepton. For the case where the recoil proton
is detected in the final state, one obtains
\begin{equation}
\frac{\dd\sigma}{\dd\Omega_p}=\frac{\vert{\AMP}\vert^2}{(4\pi)^2\,4M}\frac{\boldsymbol{p_2}^{\!2}}{\bar{ \mathfrak{D}}\vert\boldsymbol{ k_1}\vert},
\label{eq:eq13}
\end{equation}
with $\bar{ \mathfrak{D}}=(M+\epsilon_1)\vert\boldsymbol{ p_2}\vert-E_2\vert\boldsymbol{ k_1}\vert\cos\theta_p $, where $\theta_p $  is the angle between the directions 
of the lepton beam and the recoil proton, and $\dd \Omega_p$ is the differential solid angle of the scattered proton. Using the relation
\beq
\dd Q^2=\vert\boldsymbol{ k_1}\vert \vert\boldsymbol{ p_2}\vert\frac{1}{\pi}\frac{E_2+M}{\epsilon_1+M}\dd\Omega_p, 
\eeq
we obtain the following expression for the differential cross section as a function of $Q^2$
\begin{equation}
\frac{\dd\sigma}{\dd Q^2}=\frac{\vert \AMP\vert^2}{64\pi M\,\boldsymbol{k_1}^{\!2}}.
\label{eq:eq14}
\end{equation}



\subsubsection{Unpolarized cross section}



The interference of the Born diagram (\Figref{rc:elLP}, right panel) with any higher-order in $\alpha$ diagram of elastic $\ell p$ scattering (\Figref{rc:elLP}, left panel), can be expressed through a multiplicative correction, 
\begin{eqnarray} \label{eq:delta_TPE_massive}
&&\delta^{(1)}_{\rm v} = \frac{2}{ G_M^2 + \frac{\varepsilon}{\tau} G_E^2}
\bigg[ G_M \re \Delta\cG_1 + \frac{\varepsilon}{\tau} G_E \re \Delta\cG_2 \nonumber \\
&&\hspace{-0.4cm}+(1-\varepsilon_T)\left(\frac{\varepsilon_\ell}{\tau}\frac{\nu}{M^2}  G_E \re \Delta\cG_4  - G_M \re \Delta\cG_3 \right) \bigg],
\end{eqnarray}
 to the $\mathcal{O}(\al^2)$ Born cross section in \Eqref{reducedCS}.  
Here, $\re$ denotes the real part of the auxiliary amplitudes
\begin{eqnarray}\eqlab{cGamps}
 \cG_1 & = & \cG_M + \frac{\nu}{M^2} \cF_3 + \frac{m_\ell^2}{M^2} \cF_5, \label{eq:G1} \\
 \cG_2 & = & \cG_M - ( 1 + \tau ) \cF_2 + \frac{\nu}{M^2} \cF_3,  \\
 \cG_3 & = & \frac{m_\ell^2}{M^2} \cF_5 + \frac{\nu}{M^2} \cF_3,   \\
 \cG_4 & = &    \cF_4 + \frac{\nu}{M^2 (1+\tau)} \cF_5,\label{eq:G4}
\end{eqnarray}
 and
$\varepsilon_T$ is
the degree of linear polarization of transverse photons
\beq
\label{eq:transverse_epsilon}
\varepsilon_T= \frac{\nu^2  - M^4 \tau ( 1 + \tau ) ( 1 + 2 \varepsilon_\ell )}{ \nu^2 + M^4 \tau ( 1 + \tau )  ( 1 - 2 \varepsilon_\ell )}.
\eeq
In this way, one can also describe the leading TPE correction, cf.\ Eq.~\ref{delta2gTPE}. 

\subsubsection{Polarization transfer observables}\seclab{PTobservable}



The longitudinal and transverse polarization transfer asymmetries, $P_l$ and $P_t$, are defined as (with, e.g., $h=+$) 
\begin{eqnarray}
\eqlab{Pl}
P_l&=&\frac{\dd \sigma(\lambda'=+)-\dd \sigma(\lambda'=-)}{\dd \sigma(\lambda'=+)+\dd \sigma(\lambda'=-)},\\
P_t&=&\frac{\dd \sigma(S'=S_\perp)-\dd \sigma(S'=-S_\perp)}{\dd \sigma(S'=S_\perp)+\dd \sigma(S'=-S_\perp)},\quad
\eqlab{Pt}
\end{eqnarray}
where $S'=\pm S_\perp$ is the spin direction of the recoil proton in the scattering plane transverse to its momentum.
In the Born approximation, their ratio is related to
the ratio of electric to magnetic 
Sachs FFs \cite{Akhiezer:1968ek,Akhiezer:1974em}
\begin{eqnarray}
\eqlab{PTOPE}
\frac{P_t}{P_l}=-\sqrt{\frac{2\varepsilon}{\tau (1+\varepsilon)}} \frac{G_E}{G_M}.
\end{eqnarray}
Equivalent information on the proton FFs can be obtained from measuring double-spin asymmetries on a polarized proton target \cite{Dombey1969}. Taking into account the leading virtual corrections to the asymmetries, the ratio modifies, up to terms of order $\mathcal{O}(\alpha^2)$, as described in Refs.~\cite{Guichon:2003qm,Guttmann:2010au,Meziane:2010xc} as
\begin{align}
\label{PtPlradcorr}
\frac{P_t}{P_l}=-\sqrt{\frac{2\varepsilon}{\tau (1+\varepsilon)}} \Bigg[&R_{EM} +Y_E - R_{EM} Y_M \nn\\&  +\Bigg (1-\frac{2 \varepsilon}{1+\varepsilon} R_{EM} \Bigg )Y_3 \Bigg],
\end{align}
with
\begin{align}
R_{EM} \equiv \frac{G_E}{G_M},
\eqlab{REM}
\end{align}
where the shorthand notations for ratio of two-photon amplitudes relative to the magnetic FF have been used: 
\begin{align}
&Y_M \equiv \frac{\re \Delta \cG_M }{G_M}, \\ 
&Y_E \equiv \frac{\re \Delta \cG_E }{G_M}, \\ 
&Y_3 \equiv \frac{\nu}{M^2}\frac{\re \cF_3}{G_M}.
\label{eq:2gaampl}
\end{align}
Furthermore for the longitudinal polarization transfer $P_l$ separately, its expression relative to the $1 \gamma$-result $P_l^\text{Born}$ 
is given in Refs.~\cite{Guichon:2003qm,Guttmann:2010au}, up to terms of order $\mathcal{O}(\alpha^2)$, by
\begin{eqnarray}
\frac{P_l}{P_l^\text{Born}}  & =    1 - 2 \varepsilon 
\left( 1 + \frac{\varepsilon}{\tau} R_{EM}^2  \right)^{-1} 
\nn \\
&\times \left\{ 
\left[ \frac{\varepsilon}{1 + \varepsilon} \bigg( 1 - \frac{R_{EM}^2}{\tau}  \right) 
+ \frac{R_{EM}}{\tau}  \right] Y_3 
\nn \\ 
&+  \frac{R_{EM}}{\tau} \left[  Y_E  - R_{EM} Y_M \right]  \bigg\},  	\label{pl_2gamma}  
\end{eqnarray}
with
\begin{align}
& P_l^\text{Born} = \sqrt{1 - \varepsilon^2} \, 
\left( 1 + \frac{\varepsilon}{\tau} R_{EM}^2 \right)^{-1}, 
\eqlab{pl_1gamma} 
\end{align}
Note that in Eqs.~\ref{PtPlradcorr} and \ref{pl_2gamma} the lepton mass has been neglected, $m_\ell=0$, which is not a valid approximation for low-energy muon scattering \cite{Tomalak:2016kyd}.



\subsubsection{Beam and target normal single-spin asymmetries}\seclab{beam_targ_normal}
\label{SSAintro}
The target (beam) normal single-spin asymmetries (SSA) $A_n$ ($B_n$) is defined for the scattering of an unpolarized beam (a beam polarized normal to the scattering plane, with $s= \pm s_n$ the direction of incoming lepton,) off a target with normal to the scattering plane polarization $S= \pm S_n$ (an unpolarized target) 

  \begin{eqnarray}
     A_n&=&\frac{\text{d} \sigma (S=S_n)-\text{d} \sigma (S=-S_n)}{\text{d} \sigma (S=S_n)+\text{d} \sigma (S=-S_n)},\\
     B_n&=&\frac{\text{d} \sigma (s=s_n)-\text{d} \sigma (s=-s_n)}{\text{d} \sigma (s=s_n)+\text{d} \sigma (s=-s_n)}.
 \end{eqnarray}
 Both these asymmetries are zero in the  Born  approximation. Taking into account leading virtual corrections, the  asymmetries can be expressed as~\cite{Gorchtein:2004ac,Chen:2004tw,Afanasev:2005prd}
 \begin{eqnarray}
A_n&=&\sqrt{\frac{2(\varepsilon -\varepsilon_\ell)(1+\varepsilon-2\varepsilon\varepsilon_\ell)}{\tau (1-\varepsilon_\ell)^2}}
\left(1 +\frac{\varepsilon}{\tau} R_{EM}^2 \right)^{-1} \nn\\
&\times& \frac{1}{G_M} \Bigg \{- \im \bigg (\Delta \cG_2+\frac{(1+\tau)m_\ell^2}{\nu}  \cG_4\bigg ) \nn\\
&&+ R_{EM} \im \bigg(\Delta \cG_1-\frac{\tau (1+\tau) M^2}{\nu} \cF_3 \bigg ) \Bigg \},\\
B_n&=& -\frac{m_\ell}{M} \sqrt{ \frac{2 (1-\varepsilon)(\varepsilon-\varepsilon_\ell)}{ (1-\varepsilon_\ell)^2}} \frac{\sqrt{1+\tau}}{\tau} 
\left(1+\frac{\varepsilon}{\tau} R_{EM}^2 \right)^{-1} \nonumber \\
&\times& \frac{1}{G_M} \Big \{\tau \im \bigg( \cF_3+\frac{\nu}{M^2}\frac{\cF_5}{1+\tau} \bigg) + R_{EM} \im \cG_4 \Big \},\nn\\
\end{eqnarray}
 
where $\im$ denotes the imaginary part. 

